\documentclass[11pt,a4paper]{amsart}
\usepackage[T1]{fontenc}
\usepackage{lmodern}
\usepackage[brazil]{babel}
\usepackage{amsmath,amssymb,amsfonts,amsthm}
\usepackage{enumitem}
\usepackage{microtype}
\usepackage[left=2.3cm,right=2.3cm,top=2.6cm,bottom=2.6cm]{geometry}

\allowdisplaybreaks

%==================== Comandos ====================
\newcommand{\R}{\mathbb{R}}
\newcommand{\C}{\mathbb{C}}
\newcommand{\Z}{\mathbb{Z}}
\newcommand{\N}{\mathbb{N}}
\renewcommand{\d}{\mathrm{d}}
\newcommand{\norm}[1]{\lVert #1\rVert}
\newcommand{\img}{\operatorname{Im}}
\DeclareMathOperator{\diag}{diag}
\DeclareMathOperator{\tr}{tr}
\DeclareMathOperator{\posto}{posto}
\DeclareMathOperator{\Mat}{M}

%==================== Ambientes ====================
\theoremstyle{plain}
\newtheorem{teorema}{Teorema}[section]
\newtheorem{lema}[teorema]{Lema}
\newtheorem{proposicao}[teorema]{Proposição}
\newtheorem{corolario}[teorema]{Corolário}

\theoremstyle{definition}
\newtheorem{definicao}[teorema]{Definição}
\newtheorem{exemplo}[teorema]{Exemplo}
\newtheorem{exercicio}[teorema]{Exercício}

\theoremstyle{remark}
\newtheorem{observacao}[teorema]{Observação}

%==================== Título ====================
\title[Sistemas lineares de EDOs e a exponencial de matrizes]{Sistemas lineares
de equa\c{c}\~oes diferenciais ordin\'arias\\ e a exponencial de matrizes}
\author{}
\date{}

\begin{document}

\begin{abstract}
Estas notas desenvolvem a teoria dos sistemas lineares aut\^onomos
$x'(t)=Mx(t)$ a partir de um \'unico objeto central: a \emph{exponencial de
matrizes} $e^{tM}$. Definimos $e^{tM}$ por sua s\'erie de pot\^encias, provamos
suas propriedades alg\'ebricas e de deriva\c{c}\~ao, e a usamos para resolver o
problema de valor inicial. Em seguida desenvolvemos, com demonstra\c{c}\~oes
completas, a \'algebra linear necess\'aria para \emph{calcular} $e^{tM}$: os
autoespa\c{c}os generalizados, as cadeias de Jordan e a forma can\^onica de
Jordan, culminando na f\'ormula expl\'\i cita para a exponencial de um bloco de
Jordan. O texto inclui a decomposi\c{c}\~ao em cadeias --- e n\~ao apenas o
enunciado final em blocos --- e termina com exemplos trabalhados e
exerc\'\i cios.
\end{abstract}

\maketitle


\vspace{2mm}
{\small\noindent\textbf{Conven\c{c}\~ao e nota\c{c}\~ao.} $\Mat_n(\C)$ denota o
espa\c{c}o das matrizes complexas $n\times n$. Escrevemos $J_k$ para a
matriz $k\times k$ que tem $1$ em cada entrada da \emph{superdiagonal} (logo
acima da diagonal principal) e $0$ nas demais,
\[
  J_k=\begin{pmatrix}
        0 & 1 & 0 & \cdots & 0\\
        0 & 0 & 1 & \cdots & 0\\
        \vdots & & \ddots & \ddots & \vdots\\
        0 & 0 & \cdots & 0 & 1\\
        0 & 0 & \cdots & 0 & 0
      \end{pmatrix}.
\]
Note que $J_k$ \emph{n\~ao} \'e um bloco de Jordan: \'e a matriz-deslocamento
nilpotente ($J_k^{\,k}=0$). Um \emph{bloco de Jordan} de tamanho $k$ e autovalor
$\lambda$ \'e $\lambda I + J_k$; manteremos os $1$ na superdiagonal do
in\'\i cio ao fim.}

\vspace{3mm}

%=====================================================================
\section{Introdu\c{c}\~ao: por que sistemas lineares importam}
%=====================================================================

Considere o problema de valor inicial para um sistema linear aut\^onomo
\begin{equation}\label{eq:pvi}
  x'(t) = M\,x(t), \qquad x(0) = x_0,
\end{equation}
onde $M\in\Mat_n(\C)$ \'e uma matriz constante e $x\colon\R\to\C^n$. Esses
sistemas s\~ao o coment\'ario de partida de toda a teoria qualitativa de EDOs:
descrevem oscila\c{c}\~oes acopladas, circuitos, modelos populacionais
linearizados e, via lineariza\c{c}\~ao, o comportamento local de sistemas n\~ao
lineares perto de equil\'\i brios.

No caso escalar $n=1$, a equa\c{c}\~ao $x'=ax$ tem solu\c{c}\~ao
$x(t)=e^{at}x_0$. A pergunta que organiza estas notas \'e: \emph{podemos manter
exatamente essa f\'ormula no caso matricial?} A resposta \'e sim, desde que
saibamos dar sentido a $e^{tM}$. A solu\c{c}\~ao de \eqref{eq:pvi} ser\'a
\begin{equation}\label{eq:formula}
  x(t) = e^{tM}\,x_0,
\end{equation}
e as quest\~oes t\'ecnicas que precisamos responder s\~ao quatro:
\begin{enumerate}[label=(\arabic*)]
  \item O que \'e $e^{tM}$, e por que faz sentido?
  \item Por que \eqref{eq:formula} resolve o sistema, e por que \'e a
        \emph{\'unica} solu\c{c}\~ao?
  \item Como calcular $e^{tM}$ na pr\'atica?
  \item O que acontece quando $M$ \emph{n\~ao} \'e diagonaliz\'avel?
\end{enumerate}
As duas primeiras perguntas s\~ao respondidas pela teoria da exponencial
(Se\c{c}\~oes \ref{sec:exp}--\ref{sec:solucao}); a terceira e a quarta exigem
a forma can\^onica de Jordan (Se\c{c}\~oes
\ref{sec:gen}--\ref{sec:geral}). O fio condutor \'e sempre o mesmo:
\[
  \text{exponencial de matrizes}
  \;\to\;
  \text{f\'ormula da solu\c{c}\~ao}
  \;\to\;
  \text{forma de Jordan}
  \;\to\;
  \text{c\'alculo expl\'\i cito de }e^{tM}.
\]

%=====================================================================
\section{A exponencial de uma matriz}
\label{sec:exp}
%=====================================================================

\subsection{Norma de matrizes}

Fixe em $\C^n$ a norma euclidiana $\norm{x}$. A \emph{norma de operador} de
$M\in\Mat_n(\C)$ \'e
\[
  \norm{M} = \sup_{\norm{x}=1}\norm{Mx}.
\]
Ela satisfaz as propriedades que usaremos repetidamente:
\begin{equation}\label{eq:submult}
  \norm{Mx}\le \norm{M}\,\norm{x},
  \qquad
  \norm{MB}\le\norm{M}\,\norm{B},
  \qquad
  \norm{M^m}\le\norm{M}^m .
\end{equation}
A segunda desigualdade (\emph{submultiplicatividade}) segue da primeira:
\[
  \norm{MBx}\le\norm{M}\,\norm{Bx}\le\norm{M}\,\norm{B}\,\norm{x};
\]
a terceira segue por indu\c{c}\~ao. Identificando $\Mat_n(\C)$ com $\C^{n^2}$, ele \'e um
espa\c{c}o vetorial normado de dimens\~ao finita, logo \emph{completo}: toda
s\'erie absolutamente convergente de matrizes converge.

\subsection{Defini\c{c}\~ao e converg\^encia}

\begin{definicao}\label{def:exp}
Para $M\in\Mat_n(\C)$ definimos
\[
  e^{M} = \sum_{m=0}^{\infty}\frac{M^m}{m!},
  \qquad\text{e, para } t\in\R,\qquad
  e^{tM} = \sum_{m=0}^{\infty}\frac{t^m M^m}{m!},
\]
com a conven\c{c}\~ao $M^0 = I$.
\end{definicao}

\begin{proposicao}[Converg\^encia]\label{prop:conv}
Para toda $M\in\Mat_n(\C)$, a s\'erie da Defini\c{c}\~ao~\ref{def:exp}
converge absolutamente. Mais ainda, a s\'erie $\sum_m t^m M^m/m!$ converge
\emph{uniformemente} para $t$ em qualquer intervalo compacto $[-T,T]$.
\end{proposicao}

\begin{proof}
Pela estimativa \eqref{eq:submult}, para $|t|\le T$,
\[
  \norm{\frac{t^m M^m}{m!}} = \frac{|t|^m}{m!}\norm{M^m}
  \le \frac{(T\norm{M})^m}{m!}.
\]
A s\'erie num\'erica $\sum_{m\ge0}\frac{(T\norm{M})^m}{m!} = e^{T\norm{M}}$
\'e convergente. Pelo teste $M$ de Weierstrass, $\sum_m t^mM^m/m!$ converge
absoluta e uniformemente em $[-T,T]$. Tomando $T=|t|$ obtemos a
converg\^encia absoluta para cada $t$ fixo; com $t=1$, a de $e^M$.
\end{proof}

Em particular $\norm{e^{tM}}\le e^{|t|\,\norm{M}}$.

%=====================================================================
\section{Propriedades da exponencial}
\label{sec:prop}
%=====================================================================

\begin{teorema}\label{thm:propriedades}
Sejam $M,B\in\Mat_n(\C)$ e $s,t\in\R$. Valem:
\begin{enumerate}[label=\textup{(\alph*)}]
  \item $e^{0}=I$;
  \item \textup{(invari\^ancia por semelhan\c{c}a)} se $P$ \'e invert\'\i vel,
        $e^{PMP^{-1}} = P\,e^{M}\,P^{-1}$;
  \item \textup{(matrizes que comutam)} se $MB=BM$, ent\~ao
        $e^{M+B}=e^{M}e^{B}$;
  \item $e^{sM}e^{tM}=e^{(s+t)M}$;
  \item $e^{tM}$ \'e invert\'\i vel, com $\big(e^{tM}\big)^{-1}=e^{-tM}$.
\end{enumerate}
\end{teorema}

\begin{proof}
\textbf{(a)} Apenas o termo $m=0$ sobrevive: $e^0=\sum_m 0^m/m! = I$.

\medskip
\textbf{(b)} Como $(PMP^{-1})^m = P M^m P^{-1}$, a soma parcial $S_p$ satisfaz
\[
  \sum_{m=0}^{p}\frac{(PMP^{-1})^m}{m!}
  = P\left(\sum_{m=0}^{p}\frac{M^m}{m!}\right)P^{-1}.
\]
As aplica\c{c}\~oes $X\mapsto PX$ e $X\mapsto XP^{-1}$ s\~ao lineares em
dimens\~ao finita, logo cont\'\i nuas; passando ao limite $p\to\infty$ obtemos
$e^{PMP^{-1}}=Pe^{M}P^{-1}$.

\medskip
\textbf{(c)} Como ambas as s\'eries convergem absolutamente, podemos
multiplic\'a-las pelo produto de Cauchy (v\'alido em uma \'algebra de Banach):
\[
  e^{M}e^{B}
  = \left(\sum_{j\ge0}\frac{M^j}{j!}\right)\!\left(\sum_{k\ge0}\frac{B^k}{k!}\right)
  = \sum_{m\ge0}\ \sum_{j+k=m}\frac{M^j B^k}{j!\,k!}.
\]
Como $MB=BM$, vale o teorema binomial para matrizes,
\[
  (M+B)^m = \sum_{j+k=m}\binom{m}{j}M^j B^k
          = \sum_{j+k=m}\frac{m!}{j!\,k!}M^j B^k,
\]
de modo que $\sum_{j+k=m}\frac{M^jB^k}{j!\,k!} = \frac{(M+B)^m}{m!}$. Portanto
\[
  e^{M}e^{B} = \sum_{m\ge0}\frac{(M+B)^m}{m!} = e^{M+B}.
\]
\emph{A hip\'otese $MB=BM$ \'e essencial}: sem ela o binomial falha, e em geral
$e^{M+B}\neq e^Me^B$.

\medskip
\textbf{(d)} As matrizes $sM$ e $tM$ comutam; por (c),
$e^{sM}e^{tM}=e^{sM+tM}=e^{(s+t)M}$.

\medskip
\textbf{(e)} Por (d) com $s=-t$ e por (a):
$e^{tM}e^{-tM}=e^{0}=I$. Analogamente $e^{-tM}e^{tM}=I$. Logo $e^{tM}$ \'e
invert\'\i vel com inversa $e^{-tM}$.
\end{proof}

A propriedade que torna a exponencial \'util para EDOs \'e a regra de
deriva\c{c}\~ao.

\begin{teorema}[Regra de deriva\c{c}\~ao]\label{thm:derivada}
A fun\c{c}\~ao $t\mapsto e^{tM}$ \'e de classe $C^{\infty}$ e
\[
  \frac{\d}{\d t}\,e^{tM} = M\,e^{tM} = e^{tM}\,M .
\]
\end{teorema}

\begin{proof}
Consideramos a s\'erie $\sum_{m\ge0} f_m(t)$ com $f_m(t)=\dfrac{t^m M^m}{m!}$
(uma s\'erie de fun\c{c}\~oes a valores em $\Mat_n(\C)$, que podemos tratar
entrada a entrada). Cada $f_m$ \'e $C^1$ com
\[
  f_m'(t) = \frac{m\,t^{m-1}M^m}{m!} = \frac{t^{m-1}M^m}{(m-1)!}
  \qquad (m\ge1),\qquad f_0'=0.
\]
A s\'erie das derivadas \'e
\[
  \sum_{m\ge1}\frac{t^{m-1}M^m}{(m-1)!}
  = M\sum_{k\ge0}\frac{t^k M^k}{k!} = M\,e^{tM},
\]
e, pela Proposi\c{c}\~ao~\ref{prop:conv} aplicada a $M$ (a s\'erie das
derivadas \'e, a menos do fator $M$, uma exponencial), ela converge
uniformemente em cada compacto $[-T,T]$. Pelo teorema de deriva\c{c}\~ao
termo a termo (se $\sum f_m$ converge em um ponto e $\sum f_m'$ converge
uniformemente, ent\~ao $\sum f_m$ \'e deriv\'avel e
$(\sum f_m)' = \sum f_m'$), conclu\'\i mos
\[
  \frac{\d}{\d t}e^{tM} = M\,e^{tM}.
\]
Como $M$ comuta com cada soma parcial $\sum_{m\le p} t^mM^m/m!$ (pois comuta com
suas pr\'oprias pot\^encias), comuta com o limite; logo $Me^{tM}=e^{tM}M$.
Iterando, $e^{tM}$ \'e $C^\infty$.
\end{proof}

%=====================================================================
\section{Resolu\c{c}\~ao do sistema $x'=Mx$}
\label{sec:solucao}
%=====================================================================

\begin{teorema}[Solu\c{c}\~ao fundamental]\label{thm:pvi}
Para todo $x_0\in\C^n$, o problema de valor inicial \eqref{eq:pvi} tem uma
\emph{\'unica} solu\c{c}\~ao $x\colon\R\to\C^n$, dada por
\[
  x(t) = e^{tM}\,x_0.
\]
\end{teorema}

\begin{proof}
\textbf{Exist\^encia.} Seja $x(t)=e^{tM}x_0$. Pela
Teorema~\ref{thm:derivada},
\[
  x'(t) = \frac{\d}{\d t}\big(e^{tM}x_0\big) = M e^{tM}x_0 = M\,x(t),
\]
e $x(0)=e^{0}x_0 = Ix_0 = x_0$. Logo $x$ resolve \eqref{eq:pvi}.

\medskip
\textbf{Unicidade.} Seja $x(t)$ \emph{qualquer} solu\c{c}\~ao de
\eqref{eq:pvi}. Defina
\[
  y(t) = e^{-tM}\,x(t).
\]
Pela regra do produto e pela Teorema~\ref{thm:derivada} (note que
$\frac{\d}{\d t}e^{-tM} = -M e^{-tM}$),
\[
  y'(t) = -M e^{-tM}x(t) + e^{-tM}x'(t)
        = -M e^{-tM}x(t) + e^{-tM}M\,x(t).
\]
Como $M$ comuta com $e^{-tM}$ (Teorema~\ref{thm:derivada}),
$e^{-tM}M = M e^{-tM}$, de modo que os dois termos se cancelam: $y'(t)=0$ para
todo $t$. Portanto $y$ \'e constante, $y(t)=y(0)=e^{0}x_0 = x_0$. Logo
\[
  x(t) = e^{tM}y(t) = e^{tM}x_0. \qedhere
\]
\end{proof}

\begin{observacao}
A matriz $t\mapsto e^{tM}$ \'e a \emph{matriz fundamental} (ou matriz de
transi\c{c}\~ao) do sistema: suas colunas formam uma base do espa\c{c}o de
solu\c{c}\~oes, e $e^{tM}$ leva o estado inicial $x_0$ ao estado no instante
$t$. A propriedade $e^{(s+t)M}=e^{sM}e^{tM}$ \'e a lei de fluxo: evoluir por
$s+t$ \'e o mesmo que evoluir por $t$ e depois por $s$.
\end{observacao}

%=====================================================================
\section{Sistemas diagonaliz\'aveis}
\label{sec:diag}
%=====================================================================

Antes de atacar o caso geral, calculamos $e^{tM}$ nos casos simples.

\subsection{Matrizes diagonais}
Se $D=\diag(\lambda_1,\dots,\lambda_n)$, ent\~ao
$D^m=\diag(\lambda_1^m,\dots,\lambda_n^m)$, e somando entrada a entrada,
\[
  e^{tD} = \diag\!\big(e^{\lambda_1 t},\dots,e^{\lambda_n t}\big).
\]

\subsection{Matrizes diagonaliz\'aveis}
Se $M = PDP^{-1}$ com $D$ diagonal, ent\~ao pela
Teorema~\ref{thm:propriedades}(b),
\[
  e^{tM} = P\,e^{tD}\,P^{-1}
         = P\,\diag\!\big(e^{\lambda_1 t},\dots,e^{\lambda_n t}\big)\,P^{-1}.
\]
Concretamente, se $M$ possui uma base de autovetores $v_1,\dots,v_n$ com
$Mv_j=\lambda_j v_j$, ent\~ao as colunas de $P$ s\~ao os $v_j$, e a
solu\c{c}\~ao geral de $x'=Mx$ \'e a combina\c{c}\~ao linear
\[
  x(t) = \sum_{j=1}^{n} c_j\,e^{\lambda_j t}\,v_j,
  \qquad c_j\in\C.
\]
De fato, $\frac{\d}{\d t}(e^{\lambda_j t}v_j)=\lambda_j e^{\lambda_j t}v_j
= e^{\lambda_j t}Mv_j = M(e^{\lambda_j t}v_j)$, logo cada $e^{\lambda_j t}v_j$
\'e solu\c{c}\~ao, e elas formam uma base do espa\c{c}o de solu\c{c}\~oes.

Nem toda matriz \'e diagonaliz\'avel. O obst\'aculo s\~ao as matrizes
\emph{defeituosas}, em que faltam autovetores; o rem\'edio \'e a forma de
Jordan, e a pe\c{c}a que falta s\~ao as matrizes nilpotentes.

%=====================================================================
\section{Matrizes nilpotentes}
\label{sec:nilp}
%=====================================================================

\begin{definicao}
$N\in\Mat_n(\C)$ \'e \emph{nilpotente} se $N^r=0$ para algum $r\ge1$. O menor
tal $r$ \'e o \emph{\'\i ndice} de nilpot\^encia.
\end{definicao}

\begin{proposicao}\label{prop:exp-nilp}
Se $N^r=0$, ent\~ao a s\'erie de $e^{tN}$ \emph{termina}:
\[
  e^{tN} = \sum_{m=0}^{r-1}\frac{t^m N^m}{m!}
         = I + tN + \frac{t^2 N^2}{2!} + \cdots + \frac{t^{r-1}N^{r-1}}{(r-1)!}.
\]
\end{proposicao}

\begin{proof}
Para $m\ge r$ temos $N^m = N^{m-r}N^r = 0$, de modo que todos os termos a
partir de $m=r$ se anulam. Resta a soma finita.
\end{proof}

Esta \'e a observa\c{c}\~ao-chave por tr\'as dos blocos de Jordan: a
exponencial de uma parte nilpotente \'e um \emph{polin\^omio} em $t$. \'E da\'\i
que vir\~ao os fatores polinomiais nas solu\c{c}\~oes de sistemas defeituosos.

%=====================================================================
\section{Autovetores generalizados}
\label{sec:gen}
%=====================================================================

A partir de agora trabalhamos sobre $\C$, onde o polin\^omio caracter\'\i stico
$\chi_M(x)=\det(xI-M)$ se fatora completamente em fatores lineares (Teorema
Fundamental da \'Algebra). Escrevemos
\[
  \chi_M(x) = \prod_{i=1}^{k}(x-\lambda_i)^{a_i},
\]
com $\lambda_1,\dots,\lambda_k$ os autovalores \emph{distintos} e $a_i$ suas
multiplicidades alg\'ebricas, $\sum_i a_i = n$.

\begin{definicao}\label{def:gen}
Para um autovalor $\lambda$ e $m\ge1$, o \emph{$m$-\'esimo n\'ucleo iterado}
\'e $\ker(M-\lambda I)^m$. O \emph{autoespa\c{c}o generalizado} associado a
$\lambda$ \'e
\[
  K_\lambda = \ker(M-\lambda I)^{n}.
\]
\end{definicao}

\subsection{A filtra\c{c}\~ao e sua estabiliza\c{c}\~ao}

Escreva $T=M-\lambda I$. Os n\'ucleos iterados formam uma cadeia crescente
\begin{equation}\label{eq:filtracao}
  \{0\}\subseteq \ker T \subseteq \ker T^2 \subseteq \ker T^3 \subseteq\cdots,
\end{equation}
pois $T^m v=0$ implica $T^{m+1}v=T(T^m v)=0$.

\begin{lema}[Estabiliza\c{c}\~ao]\label{lema:estab}
Se $\ker T^{m}=\ker T^{m+1}$ para algum $m$, ent\~ao $\ker T^{j}=\ker T^{m}$
para todo $j\ge m$. Consequentemente a cadeia \eqref{eq:filtracao} estabiliza
em um \'\i ndice $\le n$, e $\ker T^m = \ker T^n = K_\lambda$ para todo $m\ge n$.
\end{lema}

\begin{proof}
Basta mostrar que $\ker T^{m}=\ker T^{m+1}$ implica $\ker T^{m+1}=\ker T^{m+2}$;
o resto segue por indu\c{c}\~ao. Seja $v\in\ker T^{m+2}$. Ent\~ao
$T^{m+1}(Tv)=T^{m+2}v=0$, isto \'e, $Tv\in\ker T^{m+1}=\ker T^{m}$, logo
$T^{m}(Tv)=0$, ou seja $T^{m+1}v=0$, donde $v\in\ker T^{m+1}$. A inclus\~ao
oposta j\'a vale por \eqref{eq:filtracao}.

Enquanto a cadeia cresce estritamente, a dimens\~ao aumenta de pelo menos $1$ a
cada passo; como ela \'e limitada por $n=\dim\C^n$, a cadeia deve estabilizar
em algum \'\i ndice $\le n$. A partir da\'\i\ os n\'ucleos s\~ao todos iguais a
$\ker T^n = K_\lambda$.
\end{proof}

\begin{lema}[Invari\^ancia]\label{lema:invar}
$K_\lambda$ \'e invariante por $M$: $M(K_\lambda)\subseteq K_\lambda$.
\end{lema}

\begin{proof}
$M$ comuta com $T^n=(M-\lambda I)^n$ (ambos s\~ao polin\^omios em $M$). Se
$v\in K_\lambda$, ou seja $T^n v=0$, ent\~ao $T^n(Mv)=M(T^n v)=M\cdot 0=0$,
logo $Mv\in K_\lambda$. Em particular, qualquer polin\^omio em $M$ tamb\'em
deixa $K_\lambda$ invariante.
\end{proof}

O pr\'oximo lema identifica a dimens\~ao de $K_\lambda$ com a multiplicidade
alg\'ebrica de $\lambda$. Ele depende apenas de $M-\lambda I$ (um \'unico
autovalor de cada vez), e por isso n\~ao usa a decomposi\c{c}\~ao prim\'aria ---
ser\'a, ao contr\'ario, usado para prov\'a-la.

\begin{lema}[Dimens\~ao do autoespa\c{c}o generalizado]\label{lema:dim}
Para cada autovalor $\lambda$, o operador $M-\lambda I$ \'e nilpotente em
$K_\lambda$, e
\[
  \dim K_\lambda = a_\lambda ,
\]
a multiplicidade alg\'ebrica de $\lambda$ (o expoente de $(x-\lambda)$ em
$\chi_M$).
\end{lema}

\begin{proof}
A nilpot\^encia \'e imediata: por defini\c{c}\~ao
$K_\lambda=\ker(M-\lambda I)^n$, logo $(M-\lambda I)^n=0$ em $K_\lambda$.

Para a dimens\~ao, escreva $S=M-\lambda I$ e separe $\C^n$ em n\'ucleo e imagem
iterados de $S$ (a \emph{decomposi\c{c}\~ao de Fitting}). Pela
estabiliza\c{c}\~ao (Lema~\ref{lema:estab}), $\ker S^n=\ker S^{n+1}$; por
posto-nulidade, $\dim\img S^n=\dim\img S^{n+1}$, e como
$\img S^{n+1}\subseteq\img S^n$, segue $\img S^{n+1}=\img S^n$. Assim $S$ leva
$\img S^n$ sobre si mesmo, logo $S$ \'e \emph{invert\'\i vel} em $\img S^n$.

Afirmamos que
\[
  \C^n = \underbrace{\ker S^n}_{=\,K_\lambda}\ \oplus\ \img S^n .
\]
De fato, se $v\in\ker S^n\cap\img S^n$, escreva $v=S^n w$; ent\~ao
$S^{2n}w=S^n v=0$, donde $w\in\ker S^{2n}=\ker S^n$ (estabiliza\c{c}\~ao) e
$v=S^n w=0$: a interse\c{c}\~ao \'e trivial. Como, por posto-nulidade,
$\dim\ker S^n+\dim\img S^n=n$, a soma \'e direta e preenche $\C^n$. Ambos os
somandos s\~ao invariantes por $M$ (n\'ucleo e imagem de um polin\^omio em $M$).

Finalmente, o polin\^omio caracter\'\i stico \'e multiplicativo nessa soma
direta invariante,
\[
  \chi_M = \chi_{M|_{K_\lambda}}\cdot \chi_{M|_{\img S^n}} .
\]
Como $M|_{K_\lambda}=\lambda I + S|_{K_\lambda}$ com $S|_{K_\lambda}$ nilpotente,
$\chi_{M|_{K_\lambda}}(x)=(x-\lambda)^{\dim K_\lambda}$. E como $S=M-\lambda I$
\'e invert\'\i vel em $\img S^n$, $\lambda$ n\~ao \'e autovalor de
$M|_{\img S^n}$, de modo que $(x-\lambda)\nmid\chi_{M|_{\img S^n}}$. Logo o
expoente de $(x-\lambda)$ em $\chi_M$ \'e exatamente $\dim K_\lambda$, isto \'e,
$\dim K_\lambda=a_\lambda$.
\end{proof}

\subsection{A decomposi\c{c}\~ao prim\'aria}

O resultado central desta se\c{c}\~ao \'e que $\C^n$ se decomp\~oe como soma
direta dos autoespa\c{c}os generalizados. Al\'em do Lema~\ref{lema:dim}, a
ferramenta principal \'e o lema seguinte, que diz que $M-\mu I$ ``enxerga''
apenas o autovalor $\lambda$ dentro de $K_\lambda$.

\begin{lema}\label{lema:invertivel}
Se $\mu\neq\lambda$, ent\~ao $M-\mu I$ \'e invert\'\i vel quando restrito a
$K_\lambda$. Consequentemente $(M-\mu I)^m$ \'e invert\'\i vel em $K_\lambda$
para todo $m\ge0$.
\end{lema}

\begin{proof}
Em $K_\lambda$, $T=M-\lambda I$ satisfaz $T^n=0$, ou seja, \'e nilpotente.
Escreva $M-\mu I = (M-\lambda I) - (\mu-\lambda)I = T - cI$ com
$c=\mu-\lambda\neq0$. Ent\~ao
\[
  M-\mu I = -c\Big(I - \tfrac{1}{c}T\Big),
\]
e como $\tfrac1c T$ \'e nilpotente, $I-\tfrac1c T$ \'e invert\'\i vel, com
inversa $\sum_{m=0}^{n-1}\big(\tfrac1c T\big)^m$ (s\'erie geom\'etrica finita).
Logo $M-\mu I$ \'e invert\'\i vel em $K_\lambda$; e uma pot\^encia de um operador
invert\'\i vel \'e invert\'\i vel.
\end{proof}

\begin{teorema}[Decomposi\c{c}\~ao prim\'aria]\label{thm:primaria}
$\C^n$ \'e a soma direta dos autoespa\c{c}os generalizados:
\[
  \C^n = K_{\lambda_1}\oplus K_{\lambda_2}\oplus\cdots\oplus K_{\lambda_k},
\]
e cada $K_{\lambda_i}$ \'e invariante por $M$.
\end{teorema}

A demonstra\c{c}\~ao tem duas partes independentes: primeiro mostramos que a
soma \'e \emph{direta} (os $K_{\lambda_i}$ s\~ao independentes); depois, que ela
preenche todo o $\C^n$. A invari\^ancia de cada $K_{\lambda_i}$ j\'a foi provada
no Lema~\ref{lema:invar}. A \'unica ferramenta \'e o
Lema~\ref{lema:invertivel}.

\begin{proof}[Demonstra\c{c}\~ao --- Parte 1: a soma \'e direta]
Comecemos pelo caso de dois \'\i ndices, que j\'a cont\'em a ideia.

\emph{Interse\c{c}\~ao dois a dois.} Para $i\neq j$ afirmamos
$K_{\lambda_i}\cap K_{\lambda_j}=\{0\}$. De fato, se
$v\in K_{\lambda_i}\cap K_{\lambda_j}$, ent\~ao por um lado
$v\in K_{\lambda_j}=\ker(M-\lambda_j I)^n$, isto \'e $(M-\lambda_j I)^n v=0$;
por outro lado $v\in K_{\lambda_i}$, onde $(M-\lambda_j I)^n$ \'e invert\'\i vel
(Lema~\ref{lema:invertivel}, pois $j\neq i$). Um operador invert\'\i vel s\'o
leva a $0$ o vetor $0$, logo $v=0$.

\emph{O caso geral.} Para mostrar que a soma
$K_{\lambda_1}+\cdots+K_{\lambda_k}$ \'e direta, suponha uma rela\c{c}\~ao
\begin{equation}\label{eq:relacao}
  v_1 + v_2 + \cdots + v_k = 0,
  \qquad v_m\in K_{\lambda_m},
\end{equation}
e mostremos que cada $v_m=0$. Fixe um \'\i ndice $i$ e \emph{construa a matriz}
\[
  B_i \;=\; \prod_{j\neq i}(M-\lambda_j I)^{\,n}
  \qquad(\text{um produto de pot\^encias de }M-\lambda_j I,\ j\neq i).
\]
Esses fatores s\~ao polin\^omios em $M$; em particular comutam entre si e
deixam cada $K_{\lambda_m}$ invariante (Lema~\ref{lema:invar}). A matriz $B_i$
foi feita para \emph{aniquilar todos os $K_{\lambda_j}$ exceto o $i$-\'esimo}:
\begin{itemize}[leftmargin=1.6em,topsep=2pt]
  \item Se $j\neq i$, o fator $(M-\lambda_j I)^n$ aparece em $B_i$ e anula
        $K_{\lambda_j}=\ker(M-\lambda_j I)^n$; como os fatores comutam,
        $B_i$ anula $K_{\lambda_j}$. Ou seja, $B_i v_j = 0$ para $j\neq i$.
  \item Em $K_{\lambda_i}$, cada fator $(M-\lambda_j I)^n$ ($j\neq i$) \'e
        invert\'\i vel (Lema~\ref{lema:invertivel}); um produto de operadores
        invert\'\i veis \'e invert\'\i vel, logo $B_i$ \'e invert\'\i vel em
        $K_{\lambda_i}$.
\end{itemize}
Aplicando $B_i$ \`a rela\c{c}\~ao \eqref{eq:relacao} e usando que todos os
termos com $j\neq i$ se anulam,
\[
  0 = B_i\big(v_1+\cdots+v_k\big) = \sum_{m} B_i v_m = B_i v_i .
\]
Como $v_i\in K_{\lambda_i}$ e $B_i$ \'e invert\'\i vel em $K_{\lambda_i}$, segue
$v_i=0$. O \'\i ndice $i$ era arbitr\'ario, logo todos os $v_m=0$: a soma \'e
direta. (Se os $K_{\lambda_i}$ n\~ao fossem independentes, existiria uma
rela\c{c}\~ao n\~ao trivial \eqref{eq:relacao}, e a matriz $B_i$ produziria a
contradi\c{c}\~ao $v_i=0$ para o termo sobrevivente.)
\end{proof}

\begin{proof}[Demonstra\c{c}\~ao --- Parte 2: a soma preenche $\C^n$]
Esta parte \'e agora imediata, por contagem de dimens\~oes. Pela Parte 1 a
soma \'e direta, logo sua dimens\~ao \'e a soma das dimens\~oes dos somandos.
Usando o Lema~\ref{lema:dim} (cada $\dim K_{\lambda_i}=a_i$) e que a soma das
multiplicidades alg\'ebricas \'e o grau de $\chi_M$, isto \'e $n$:
\[
  \dim\big(K_{\lambda_1}\oplus\cdots\oplus K_{\lambda_k}\big)
  = \sum_{i=1}^{k}\dim K_{\lambda_i}
  = \sum_{i=1}^{k} a_i
  = \deg\chi_M = n .
\]
Um subespa\c{c}o de $\C^n$ com dimens\~ao $n$ \'e o pr\'oprio $\C^n$. Portanto
$K_{\lambda_1}\oplus\cdots\oplus K_{\lambda_k}=\C^n$.
\end{proof}

%=====================================================================
\section{Cadeias de Jordan}
\label{sec:cadeias}
%=====================================================================

Pela decomposi\c{c}\~ao prim\'aria, o estudo de $M$ se reduz, em cada
$K_{\lambda_i}$, ao do operador \emph{nilpotente} $T_i=M-\lambda_i I$. O
teorema seguinte \'e o cora\c{c}\~ao da forma de Jordan: todo operador
nilpotente admite uma base feita de \emph{cadeias}.

\begin{definicao}\label{def:cadeia}
Seja $N$ nilpotente em um espa\c{c}o $V$. Uma \emph{cadeia de Jordan} de
comprimento $d$ gerada por $v$ \'e a lista
\[
  v,\ Nv,\ N^2 v,\ \dots,\ N^{d-1}v,
  \qquad\text{com } N^{d-1}v\neq 0,\ N^{d}v=0.
\]
O vetor $N^{d-1}v\in\ker N$ \'e o \emph{fundo} da cadeia (um autovetor); $v$ \'e
o \emph{topo}.
\end{definicao}

\begin{teorema}[Base de cadeias para operadores nilpotentes]\label{thm:nilp-jordan}
Seja $N$ um operador nilpotente em um espa\c{c}o vetorial $V$ de dimens\~ao
finita. Ent\~ao $V$ admite uma base formada por uni\~ao de cadeias de Jordan
de $N$.
\end{teorema}

\paragraph{A ideia da demonstra\c{c}\~ao.}
Queremos organizar uma base de $V$ em ``colunas'' --- as cadeias --- em que
$N$ desce um degrau de cada vez,
\[
  \underbrace{v}_{\text{topo}} \;\xrightarrow{\,N\,}\; Nv
  \;\xrightarrow{\,N\,}\; \cdots
  \;\xrightarrow{\,N\,}\; \underbrace{N^{d-1}v}_{\text{fundo}\,\in\,\ker N}
  \;\xrightarrow{\,N\,}\; 0 .
\]
A demonstra\c{c}\~ao \'e por indu\c{c}\~ao, e a chave \'e olhar para
$U=\img N$: aplicar $N$ ``empurra tudo um degrau para baixo'', de modo que
$U$ \'e menor que $V$ e nele a hip\'otese de indu\c{c}\~ao j\'a fornece cadeias.
A partir dessas cadeias de $U$ fazemos duas coisas: \emph{levantamos} cada
uma um degrau (achando uma pr\'e-imagem do seu topo, o que prolonga a cadeia
para cima) e depois \emph{completamos} o que falta no n\'ucleo $\ker N$ com
cadeias de comprimento $1$. O cuidado todo est\'a em mostrar que o resultado
\'e de fato uma base --- e \'e a\'\i\ que a contagem por posto-nulidade entra.

\begin{proof}
Indu\c{c}\~ao em $\dim V$. Se $\dim V=0$ nada h\'a a provar (base vazia).
Suponha o resultado v\'alido para todo espa\c{c}o de dimens\~ao $<\dim V$.

\emph{Reduzir a $U=\img N$.} Seja $U=\img N = N(V)$. Como $N$ \'e nilpotente em
$V\neq\{0\}$, ela n\~ao \'e sobrejetora (uma matriz nilpotente \emph{e}
invert\'\i vel for\c{c}aria $V=\{0\}$), logo $\dim U<\dim V$. Al\'em disso $U$
\'e invariante por $N$ e $N|_U$ \'e nilpotente. Pela hip\'otese de
indu\c{c}\~ao, $U$ admite uma base formada por cadeias; sejam elas geradas por
$w_1,\dots,w_p\in U$, com comprimentos $k_1,\dots,k_p$:
\begin{equation}\label{eq:base-U}
  \big\{\,N^{j}w_i : 1\le i\le p,\ 0\le j\le k_i-1\,\big\}
  \quad\text{\'e base de } U,
  \qquad N^{k_i}w_i=0.
\end{equation}

\emph{Passo 1: levantar cada cadeia um degrau.} Cada topo $w_i$ pertence a
$U=\img N$, logo \emph{tem uma pr\'e-imagem}: existe $v_i\in V$ com $Nv_i=w_i$.
Trocando o topo $w_i$ por $v_i$, a cadeia gerada por $v_i$ \'e a de $w_i$
prolongada um degrau para cima:
\[
  v_i,\ Nv_i=w_i,\ N^2 v_i = N w_i,\ \dots,\ N^{k_i}v_i = N^{k_i-1}w_i,
  \qquad N^{k_i+1}v_i = N^{k_i}w_i = 0,
\]
agora de comprimento $k_i+1$. Os fundos dessas cadeias prolongadas s\~ao
\[
  b_i := N^{k_i}v_i = N^{k_i-1}w_i \in \ker N ,
\]
e eles fazem parte da base \eqref{eq:base-U} de $U$ (s\~ao os fundos das
cadeias originais), portanto s\~ao linearmente independentes.

\emph{Passo 2: completar o n\'ucleo.} Os $b_1,\dots,b_p$ s\~ao vetores
independentes de $\ker N$, mas podem n\~ao ger\'a-lo. Estenda-os a uma base
de $\ker N$ acrescentando vetores $z_1,\dots,z_q\in\ker N$. Cada $z_l$, sozinho,
\'e uma cadeia de comprimento $1$ (seu pr\'oprio fundo).

\emph{Afirma\c{c}\~ao.} A uni\~ao das cadeias prolongadas com as cadeias
triviais, isto \'e
\[
  \mathcal{B} = \big\{N^{j}v_i : 1\le i\le p,\ 0\le j\le k_i\big\}
                \cup \{z_1,\dots,z_q\},
\]
\'e uma base de $V$. (Note que $\mathcal{B}$ \'e exatamente uma uni\~ao de
cadeias de $N$, como quer\'\i amos.)

\emph{Contagem.} O n\'umero de vetores \'e
$|\mathcal{B}| = \sum_i(k_i+1)+q = \big(\sum_i k_i\big)+p+q$. Ora,
$\dim U=\sum_i k_i$ (tamanho da base \eqref{eq:base-U}) e $\dim\ker N=p+q$
(base estendida no Passo 2). Pela rela\c{c}\~ao posto-nulidade,
\[
  \dim V=\dim\img N+\dim\ker N = \textstyle\sum_i k_i + (p+q) = |\mathcal{B}|.
\]
Como $|\mathcal{B}|=\dim V$, basta provar que $\mathcal{B}$ \'e linearmente
independente para concluir que \'e base.

\emph{Independ\^encia.} Suponha uma combina\c{c}\~ao linear nula
\begin{equation}\label{eq:comb-nula}
  \sum_{i=1}^{p}\sum_{j=0}^{k_i} a_{ij}\,N^{j}v_i
  + \sum_{l=1}^{q} c_l\,z_l = 0.
\end{equation}
A estrat\'egia \'e aplicar $N$ para ``empurrar tudo para baixo'' e usar a
base de $U$. Aplicando $N$ a \eqref{eq:comb-nula}: os termos $z_l$ desaparecem
(est\~ao em $\ker N$), e $N(N^{j}v_i)=N^{j+1}v_i$, com $N^{k_i+1}v_i=0$
matando o termo de topo $j=k_i$. Sobra
\[
  \sum_{i=1}^{p}\sum_{j=0}^{k_i-1} a_{ij}\,N^{j+1}v_i = 0,
  \qquad\text{ou seja}\qquad
  \sum_{i=1}^{p}\sum_{j=0}^{k_i-1} a_{ij}\,N^{j}w_i = 0,
\]
pois $N^{j+1}v_i=N^{j}w_i$. Esses s\~ao precisamente os vetores da base
\eqref{eq:base-U} de $U$; logo todos os coeficientes se anulam:
\[
  a_{ij}=0 \quad\text{para } 0\le j\le k_i-1 .
\]
Voltando a \eqref{eq:comb-nula}, sobrevivem apenas os termos de topo $j=k_i$
e os $z_l$:
\[
  \sum_{i=1}^{p} a_{i,k_i}\,N^{k_i}v_i + \sum_{l=1}^{q} c_l z_l
  = \sum_{i=1}^{p} a_{i,k_i}\,b_i + \sum_{l=1}^{q} c_l z_l = 0.
\]
Mas $\{b_1,\dots,b_p,z_1,\dots,z_q\}$ \'e base de $\ker N$ (Passo 2), portanto
$a_{i,k_i}=0$ e $c_l=0$ para todos $i,l$. Assim todos os coeficientes de
\eqref{eq:comb-nula} s\~ao nulos: $\mathcal{B}$ \'e independente, logo \'e base
de $V$. Isto conclui a indu\c{c}\~ao.
\end{proof}

\begin{observacao}[A cadeia produz um bloco de Jordan]\label{obs:bloco}
Seja $v,Nv,\dots,N^{r-1}v$ uma cadeia de comprimento $r$. Ordene-a do fundo
para o topo:
\[
  w_1 = N^{r-1}v,\quad w_2 = N^{r-2}v,\quad\dots,\quad w_r = v
  \qquad (\text{isto \'e, } w_s = N^{r-s}v).
\]
Ent\~ao $Nw_1 = N^r v = 0$ e, para $s\ge2$, $Nw_s = N^{r-s+1}v = w_{s-1}$. Na
base ordenada $(w_1,\dots,w_r)$, a matriz de $N$ \'e precisamente a
matriz-deslocamento $J_r$ (os $1$ na superdiagonal),
\[
  [N]
  = J_r
  = \begin{pmatrix}
      0 & 1 & 0 & \cdots & 0\\
      0 & 0 & 1 & \cdots & 0\\
      \vdots & & \ddots & \ddots & \vdots\\
      0 & 0 & \cdots & 0 & 1\\
      0 & 0 & \cdots & 0 & 0
    \end{pmatrix}.
\]
\end{observacao}

%=====================================================================
\section{A forma can\^onica de Jordan}
\label{sec:jordan}
%=====================================================================

Reunindo a decomposi\c{c}\~ao prim\'aria com a base de cadeias, obtemos o
teorema de estrutura.

\begin{teorema}[Forma can\^onica de Jordan]\label{thm:jordan}
Para toda $M\in\Mat_n(\C)$ existe uma matriz invert\'\i vel $P$ tal que
\[
  P^{-1}MP = \mathcal{J} = \diag\big(\lambda_1 I + J_{k_1},\ \dots,\
  \lambda_s I + J_{k_s}\big),
\]
isto \'e, $\mathcal{J}$ \'e diagonal em blocos, cada bloco sendo um bloco de
Jordan
\[
  \lambda I + J_k =
  \begin{pmatrix}
    \lambda & 1 & & \\
    & \lambda & \ddots & \\
    & & \ddots & 1\\
    & & & \lambda
  \end{pmatrix}
  \qquad(k\times k),
\]
com $\lambda$ um autovalor de $M$ e $J_k$ a matriz-deslocamento ($J_k^{\,k}=0$).
\end{teorema}

\begin{proof}
Pela Teorema~\ref{thm:primaria},
$\C^n=\bigoplus_{i=1}^{k}K_{\lambda_i}$, com cada $K_{\lambda_i}$ invariante por
$M$. Em $K_{\lambda_i}$, o operador $T_i=M-\lambda_i I$ \'e nilpotente
(Lema~\ref{lema:dim}); pela
Teorema~\ref{thm:nilp-jordan}, $K_{\lambda_i}$ tem uma base feita de
cadeias de $T_i$. Ordenando cada cadeia (de comprimento $k$) do fundo para o
topo (Observa\c{c}\~ao~\ref{obs:bloco}), a matriz de $T_i$ nessa cadeia \'e a
matriz-deslocamento $J_k$; como $M=\lambda_i I + T_i$, a matriz de $M$ na
cadeia \'e o bloco de Jordan $\lambda_i I + J_k$.

Justapondo as bases de cadeias de todos os $K_{\lambda_i}$ obtemos uma base de
$\C^n$; seja $P$ a matriz cujas colunas s\~ao esses vetores de base. Nessa
base, $M$ \'e diagonal em blocos, com um bloco de Jordan por cadeia. Isto \'e
$P^{-1}MP=\mathcal{J}$.
\end{proof}

\begin{observacao}[Diciton\'ario das multiplicidades]
Para um autovalor $\lambda$: o n\'umero de blocos de Jordan de $\lambda$ \'e a
\emph{multiplicidade geom\'etrica} $\dim\ker(M-\lambda I)$ (cada cadeia
contribui exatamente um autovetor, seu fundo); a soma dos tamanhos dos blocos
de $\lambda$ \'e a \emph{multiplicidade alg\'ebrica} $a_\lambda=\dim
K_\lambda$. A matriz \'e diagonaliz\'avel se, e somente se, todos os blocos
t\^em tamanho $1$, ou seja, multiplicidade geom\'etrica $=$ alg\'ebrica para
todo $\lambda$.
\end{observacao}

%=====================================================================
\section{A exponencial de um bloco de Jordan}
\label{sec:exp-jordan}
%=====================================================================

Agora colhemos o fruto: calcular explicitamente a exponencial de um bloco de
Jordan.

\begin{teorema}\label{thm:exp-bloco}
Seja $\lambda I + J_r$ um bloco de Jordan $r\times r$, onde $J_r$ \'e a
matriz-deslocamento ($J_r^{\,r}=0$). Ent\~ao
\[
  e^{t(\lambda I + J_r)} = e^{\lambda t}\sum_{m=0}^{r-1}\frac{t^m J_r^{\,m}}{m!}
         = e^{\lambda t}
           \begin{pmatrix}
             1 & t & \dfrac{t^2}{2!} & \cdots & \dfrac{t^{r-1}}{(r-1)!}\\[2mm]
             0 & 1 & t & \cdots & \dfrac{t^{r-2}}{(r-2)!}\\[1mm]
             0 & 0 & 1 & \cdots & \dfrac{t^{r-3}}{(r-3)!}\\
             \vdots & & & \ddots & \vdots\\
             0 & 0 & 0 & \cdots & 1
           \end{pmatrix}.
\]
\end{teorema}

\begin{proof}
As matrizes $\lambda I$ e $J_r$ comutam, logo pela
Teorema~\ref{thm:propriedades}(c),
\[
  e^{t(\lambda I + J_r)} = e^{t\lambda I + tJ_r} = e^{t\lambda I}\,e^{tJ_r}
         = e^{\lambda t}\,e^{tJ_r}.
\]
Como $J_r^{\,r}=0$, a Proposi\c{c}\~ao~\ref{prop:exp-nilp} d\'a
$e^{tJ_r}=\sum_{m=0}^{r-1}t^m J_r^{\,m}/m!$. Resta identificar as entradas. A
pot\^encia $J_r^{\,m}$ \'e a matriz com $1$ na $m$-\'esima superdiagonal e $0$
nas demais (deslocar $m$ vezes), isto \'e $(J_r^{\,m})_{ij}=1$ se $j-i=m$ e $0$
caso contr\'ario. Portanto
\[
  (e^{tJ_r})_{ij}=\sum_{m\ge0}\frac{t^m}{m!}(J_r^{\,m})_{ij}
  = \begin{cases}\dfrac{t^{\,j-i}}{(j-i)!}, & j\ge i,\\[2mm] 0, & j<i,\end{cases}
\]
que \'e exatamente a matriz triangular superior do enunciado, multiplicada por
$e^{\lambda t}$.
\end{proof}

\begin{corolario}\label{cor:fatores}
As solu\c{c}\~oes de $x'=Mx$ associadas a uma cadeia de Jordan de tamanho $r$
para o autovalor $\lambda$ s\~ao combina\c{c}\~oes de termos
\[
  e^{\lambda t}\,p(t),
\]
onde $p$ \'e um polin\^omio vetorial de grau $\le r-1$. Em particular, fatores
polinomiais aparecem \emph{exatamente} quando o bloco tem tamanho $\ge2$, isto
\'e, quando $M$ \'e defeituosa.
\end{corolario}

%=====================================================================
\section{A f\'ormula geral para $e^{tM}$}
\label{sec:geral}
%=====================================================================

\begin{teorema}\label{thm:geral}
Seja $M=P\,\mathcal{J}\,P^{-1}$ a forma de Jordan, com
$\mathcal{J}=\diag\big(\lambda_1 I + J_{k_1},\dots,\lambda_s I + J_{k_s}\big)$.
Ent\~ao
\[
  e^{tM} = P\,e^{t\mathcal{J}}\,P^{-1},
  \qquad
  e^{t\mathcal{J}} = \diag\!\Big(e^{t(\lambda_1 I + J_{k_1})},\dots,
  e^{t(\lambda_s I + J_{k_s})}\Big),
\]
e cada bloco $e^{t(\lambda_b I + J_{k_b})}$ \'e dado pela
Teorema~\ref{thm:exp-bloco}. A solu\c{c}\~ao de $x'=Mx,\ x(0)=x_0$, \'e
$x(t)=P\,e^{t\mathcal{J}}\,P^{-1}x_0$.
\end{teorema}

\begin{proof}
A f\'ormula $e^{tM}=P\,e^{t\mathcal{J}}\,P^{-1}$ \'e a
Teorema~\ref{thm:propriedades}(b). Para uma matriz diagonal em blocos, a
$m$-\'esima pot\^encia \'e diagonal nos mesmos blocos, cada um elevado a $m$;
somando entrada a entrada a s\'erie da exponencial, $e^{t\mathcal{J}}$ \'e
diagonal em blocos com os $e^{t(\lambda_b I + J_{k_b})}$ na diagonal. O restante
\'e a Teorema~\ref{thm:pvi}.
\end{proof}

Isto fornece um m\'etodo completo: \emph{(i)} ache autovalores e autoespa\c{c}os
generalizados; \emph{(ii)} em cada um, construa cadeias de Jordan e monte $P$;
\emph{(iii)} escreva $e^{t\mathcal{J}}$ bloco a bloco; \emph{(iv)} retorne via
$e^{tM}=P\,e^{t\mathcal{J}}\,P^{-1}$.

%=====================================================================
\section{Sistemas reais e autovalores complexos}
\label{sec:real}
%=====================================================================

Se $M\in\Mat_n(\R)$, ainda assim a forma de Jordan vive naturalmente sobre
$\C$. Mas as solu\c{c}\~oes de interesse f\'\i sico s\~ao reais; eis como
recuper\'a-las.

\begin{itemize}[leftmargin=1.4em]
  \item Autovalores complexos de uma matriz real ocorrem em \emph{pares
        conjugados}: se $Mv=\lambda v$ com $\lambda=\alpha+i\beta$, ent\~ao,
        conjugando, $M\bar v=\bar\lambda\bar v$.
  \item Dada uma solu\c{c}\~ao complexa $z(t)=e^{\lambda t}v$ de $x'=Mx$ com $M$
        real, suas partes real e imagin\'aria, $\operatorname{Re}z$ e
        $\operatorname{Im}z$, s\~ao solu\c{c}\~oes \emph{reais}. De fato,
        $z'=Mz$ com $M$ real implica $(\operatorname{Re}z)'=M\operatorname{Re}z$
        e $(\operatorname{Im}z)'=M\operatorname{Im}z$.
  \item Para um par $\alpha\pm i\beta$, as solu\c{c}\~oes reais envolvem
        \[
          e^{\alpha t}\cos(\beta t), \qquad e^{\alpha t}\sin(\beta t),
        \]
        produzindo espirais (se $\alpha\neq0$) ou oscila\c{c}\~oes puras (se
        $\alpha=0$).
\end{itemize}
Existe uma \emph{forma can\^onica real} que substitui cada par de blocos
complexos conjugados por um bloco real com mini-blocos
$\big(\begin{smallmatrix}\alpha&-\beta\\ \beta&\alpha\end{smallmatrix}\big)$;
n\~ao a desenvolveremos aqui para n\~ao desviar do c\'alculo central de
$e^{tM}$. O Exemplo~\ref{ex:rotacao} ilustra o caso $\alpha=0$.

%=====================================================================
\section{Exemplos}
\label{sec:exemplos}
%=====================================================================

\begin{exemplo}[Matriz diagonal]
Seja $M=\begin{pmatrix}2&0\\0&-1\end{pmatrix}$. Ent\~ao
\[
  e^{tM}=\begin{pmatrix}e^{2t}&0\\0&e^{-t}\end{pmatrix},
\]
e a solu\c{c}\~ao de $x'=Mx$ \'e
$x(t)=\big(x_1(0)e^{2t},\,x_2(0)e^{-t}\big)$. A primeira coordenada cresce, a
segunda decai: a origem \'e um ponto de sela.
\end{exemplo}

\begin{exemplo}[Diagonaliz\'avel, mas n\~ao diagonal]\label{ex:diag}
Seja $M=\begin{pmatrix}1&1\\4&1\end{pmatrix}$. O polin\^omio
caracter\'\i stico \'e $\lambda^2-2\lambda-3=(\lambda-3)(\lambda+1)$, com
autovalores $3$ e $-1$. Resolvendo $(M-3I)v=0$ e $(M+I)v=0$,
\[
  v_3=\begin{pmatrix}1\\2\end{pmatrix},\qquad
  v_{-1}=\begin{pmatrix}1\\-2\end{pmatrix}.
\]
Tome $P=\begin{pmatrix}1&1\\2&-2\end{pmatrix}$, $\det P=-4$,
$P^{-1}=\begin{pmatrix}1/2&1/4\\ 1/2&-1/4\end{pmatrix}$. Ent\~ao
$e^{tM}=P\diag(e^{3t},e^{-t})P^{-1}$, isto \'e
\[
  e^{tM}=\begin{pmatrix}
    \dfrac{e^{3t}+e^{-t}}{2} & \dfrac{e^{3t}-e^{-t}}{4}\\[2mm]
    e^{3t}-e^{-t} & \dfrac{e^{3t}+e^{-t}}{2}
  \end{pmatrix}.
\]
\emph{Verifica\c{c}\~ao:} em $t=0$ d\'a $I$; derivando em $t=0$ recupera-se
$M$.
\end{exemplo}

\begin{exemplo}[Um bloco de Jordan $2\times2$]
Seja $M=\begin{pmatrix}\lambda&1\\0&\lambda\end{pmatrix}=\lambda I+J_2$ com
$J_2=\big(\begin{smallmatrix}0&1\\0&0\end{smallmatrix}\big)$, $J_2^{\,2}=0$.
Ent\~ao
\[
  e^{tM}=e^{\lambda t}(I+tJ_2)
        =e^{\lambda t}\begin{pmatrix}1&t\\0&1\end{pmatrix}.
\]
A solu\c{c}\~ao $x(t)=e^{tM}x_0$ cont\'em o fator polinomial $t$, sinal da
defici\^encia de autovetores.
\end{exemplo}

\begin{exemplo}[Um bloco de Jordan $3\times3$]
Seja $M=\begin{pmatrix}\lambda&1&0\\0&\lambda&1\\0&0&\lambda\end{pmatrix}$.
Ent\~ao
\[
  e^{tM}=e^{\lambda t}
  \begin{pmatrix}1&t&t^2/2\\0&1&t\\0&0&1\end{pmatrix}.
\]
\end{exemplo}

\begin{exemplo}[Encontrando uma cadeia: caso $3\times3$ triangular]
\label{ex:cadeia3}
Seja $M=\begin{pmatrix}1&1&0\\0&1&1\\0&0&1\end{pmatrix}$. O \'unico autovalor
\'e $\lambda=1$, e $M-I=\begin{pmatrix}0&1&0\\0&0&1\\0&0&0\end{pmatrix}=J_3$, com
$(M-I)^2=\begin{pmatrix}0&0&1\\0&0&0\\0&0&0\end{pmatrix}$ e $(M-I)^3=0$.
Procuramos $v$ com $(M-I)^2 v\neq0$: tome $v=e_3=(0,0,1)^\top$, pois
$(M-I)^2 e_3 = e_1\neq0$ e $(M-I)^3 e_3=0$. A cadeia, do fundo ao topo, \'e
\[
  w_1=(M-I)^2 e_3 = e_1,\quad w_2=(M-I)e_3=e_2,\quad w_3=e_3,
\]
que j\'a \'e a base can\^onica: $M$ j\'a est\'a em forma de Jordan, e
$e^{tM}=e^{t}\big(\begin{smallmatrix}1&t&t^2/2\\0&1&t\\0&0&1\end{smallmatrix}\big)$.
\end{exemplo}

\begin{exemplo}[Encontrando uma cadeia: caso defeituoso n\~ao triangular]
\label{ex:cadeia2}
Seja $M=\begin{pmatrix}3&1\\-1&1\end{pmatrix}$. O polin\^omio
caracter\'\i stico \'e $\lambda^2-4\lambda+4=(\lambda-2)^2$: autovalor $2$ com
multiplicidade alg\'ebrica $2$. Mas
$M-2I=\begin{pmatrix}1&1\\-1&-1\end{pmatrix}$ tem posto $1$, logo um \'unico
autovetor (a menos de escala): $v_1=(1,-1)^\top$. A matriz \'e defeituosa.
Buscamos um autovetor generalizado $v_2$ com $(M-2I)v_2=v_1$:
\[
  \begin{pmatrix}1&1\\-1&-1\end{pmatrix}v_2=\begin{pmatrix}1\\-1\end{pmatrix}
  \ \Longrightarrow\ v_2=\begin{pmatrix}1\\0\end{pmatrix}\ \text{(uma escolha).}
\]
A cadeia \'e $w_1=v_1=(1,-1)^\top$ (fundo), $w_2=v_2=(1,0)^\top$ (topo). Com
$P=\begin{pmatrix}1&1\\-1&0\end{pmatrix}$, $\det P=1$,
$P^{-1}=\begin{pmatrix}0&-1\\1&1\end{pmatrix}$, obtemos
$P^{-1}MP=\begin{pmatrix}2&1\\0&2\end{pmatrix}$. Logo
\[
  e^{tM}=P\,e^{2t}\begin{pmatrix}1&t\\0&1\end{pmatrix}P^{-1}
        =e^{2t}\begin{pmatrix}1+t & t\\ -t & 1-t\end{pmatrix}.
\]
\emph{Verifica\c{c}\~ao:} em $t=0$ d\'a $I$; a derivada em $t=0$ \'e
$\big(\begin{smallmatrix}3&1\\-1&1\end{smallmatrix}\big)=M$.
\end{exemplo}

\begin{exemplo}[Matriz real com autovalores complexos]\label{ex:rotacao}
Seja $M=\begin{pmatrix}0&-1\\1&0\end{pmatrix}$. Como
$M^2=-I$, temos $M^{2k}=(-1)^k I$ e $M^{2k+1}=(-1)^k M$. Logo
\[
  e^{tM}=\sum_{m\ge0}\frac{t^m M^m}{m!}
  =\Big(\sum_{k\ge0}\frac{(-1)^k t^{2k}}{(2k)!}\Big)I
  +\Big(\sum_{k\ge0}\frac{(-1)^k t^{2k+1}}{(2k+1)!}\Big)M
  =\cos t\,I+\sin t\,M,
\]
isto \'e
\[
  e^{tM}=\begin{pmatrix}\cos t & -\sin t\\ \sin t & \cos t\end{pmatrix}.
\]
Os autovalores s\~ao $\pm i$; as solu\c{c}\~oes giram em c\'\i rculos em torno
da origem (um centro).
\end{exemplo}

%=====================================================================
\section{Exerc\'\i cios}
\label{sec:exercicios}
%=====================================================================

\subsection*{B\'asicos}
\begin{exercicio}
Mostre que $e^{tM}$ comuta com $M$ para todo $t$.
\end{exercicio}
\begin{exercicio}
Prove diretamente, a partir da defini\c{c}\~ao, que
$e^{sM}e^{tM}=e^{(s+t)M}$.
\end{exercicio}
\begin{exercicio}
Calcule $e^{tM}$ para $M=\diag(\lambda_1,\dots,\lambda_n)$.
\end{exercicio}
\begin{exercicio}
Calcule $e^{tN}$ para $N=\big(\begin{smallmatrix}0&1\\0&0\end{smallmatrix}\big)$.
\end{exercicio}

\subsection*{Intermedi\'arios}
\begin{exercicio}
Diagonalize $M=\big(\begin{smallmatrix}0&1\\-2&-3\end{smallmatrix}\big)$ e
resolva $x'=Mx$.
\end{exercicio}
\begin{exercicio}
Calcule $e^{tM}$ para o bloco $M=\big(\begin{smallmatrix}\lambda&1\\0&\lambda
\end{smallmatrix}\big)$ e para o bloco $3\times3$ correspondente.
\end{exercicio}
\begin{exercicio}
Encontre uma cadeia de Jordan para $M=\big(\begin{smallmatrix}3&1\\-1&1
\end{smallmatrix}\big)$ (compare com o Exemplo~\ref{ex:cadeia2}) e para uma
matriz $3\times3$ com um \'unico autovalor e um \'unico autovetor.
\end{exercicio}

%=====================================================================
\section*{Refer\^encias}
%=====================================================================
\begin{itemize}[leftmargin=1.4em]
  \item G.~Teschl, \emph{Ordinary Differential Equations and Dynamical
        Systems}, cap.~3 e ap\^endice sobre forma de Jordan --- refer\^encia
        unificada mais pr\'oxima destas notas (exponencial de matrizes,
        sistemas aut\^onomos, autovetores generalizados e cadeias).
  \item W.~Boyce, R.~DiPrima, \emph{Elementary Differential Equations and
        Boundary Value Problems}, se\c{c}\~ao sobre sistemas lineares com
        autovalores repetidos --- exemplos e motiva\c{c}\~ao amig\'aveis.
  \item S.~Weintraub, \emph{Jordan Canonical Form: Application to Differential
        Equations} --- exemplos adicionais conectando blocos de Jordan a
        sistemas de EDO.
\end{itemize}

\end{document}
